\BourbakiExerciseGroup

\begin{enumerate}
\def\labelenumi{\arabic{enumi})}
\item
  Let X be a topological space, S a subbase of the topology of X (§ 2, no. 3). Show that X is quasi-compact if every open covering of X formed by sets belonging to S contains a finite covering of X. (Observe that if an ultrafilter U on X were not convergent, then every \(x \in X\) would belong to a set of S not belonging to U; and use the corollary to Proposition 5 of § 6).
\item
  Show that a well-ordered set X with the topology \(\mathcal{C}_{-}(\mathbf{X})\) (§ 2, Exercise 5) is locally compact, and is compact if and only if X has a greatest element. Deduce that every set can be topologized in such a way that it becomes a compact space.
\item
  \begin{enumerate}
  \def\labelenumii{\alph{enumii})}
  \tightlist
  \item
    Let X be a Hausdorff space, and let A, B be two disjoint compact subsets of X. Show that there is a neighbourhood of A and a neighbourhood of B which have no common point.
  \end{enumerate}
\end{enumerate}

\begin{enumerate}
\def\labelenumi{\alph{enumi})}
\setcounter{enumi}{1}
\tightlist
\item
  Let X be a regular space and A a compact subset of X, and let B be a closed subset of X which does not meet A. Show that there is a neighbourhood of A and a neighbourhood of B which have no common point.
\end{enumerate}

\begin{enumerate}
\def\labelenumi{\arabic{enumi})}
\setcounter{enumi}{3}
\item
  Show that in the space X defined in Exercise 6 of § 7, which is regular (§ 8, Exercise 16), every compact subset is finite.
\item
  *a) Show that the non-Hausdorff accessible space Y = X/S defined in Exercise 7 of § 8 is such that every point of Y has a compact neighbourhood; that the images α, β of 1 and -1 in Y have compact neighbourhoods which are not closed; that the intersection of a compact neighbourhood of α and a compact neighbourhood of β is not quasi-compact, and that the union of these two neighbourhoods is not compact.*
\end{enumerate}

\begin{enumerate}
\def\labelenumi{\alph{enumi})}
\setcounter{enumi}{1}
\item
  Let X be the sum of N and an infinite set A. For each \(x \in X\) let \(\mathfrak{S}(x)\) be the filter base defined as follows: if \(x \in N\) , then \(\mathfrak{S}(x)\) consists of the single set \(\{x\}\) ; if \(x \in A\) and if \(S_{x,n}\) denotes the set whose elements are x and the integers \(\geqslant n\) , then \(\mathfrak{S}(x)\) consists of the sets \(S_{x,n}\) where n runs through N. Show that there is a topology on X for which \(\mathfrak{S}(x)\) is a fundamental system of neighbourhoods of x for every \(x \in X\) , and that X is accessible but not Hausdorff in this topology. Show also that X is not quasi-compact, but that X has dense compact subsets.
\item
  Show that two topologies on a set X, with respect to each of which X is accessible and quasi-compact, can be distinct but comparable {[}cf.~§ 8, Exercises 7 and 5 b){]}.
\end{enumerate}

\begin{enumerate}
\def\labelenumi{\arabic{enumi})}
\setcounter{enumi}{5}
\item
  Let X, Y be two topological spaces, and let A (resp. B) be a quasi-compact subset of X (resp. Y). Show that if U is any neighbourhood of \(A \times B\) in \(X \times Y\) then there is a neighbourhood V of A in X and a neighbourhood W of B in Y such that \(V \times W \subset U\) .
\item
  \begin{enumerate}
  \def\labelenumii{\alph{enumii})}
  \tightlist
  \item
    Give an example of an inverse system \((\mathbf{X}_{\alpha}, f_{\alpha\beta})\) of quasi-compact spaces whose inverse limit is empty (observe that in the space of Example 2 of no. 1, every subspace is quasi-compact).
  \end{enumerate}
\end{enumerate}

\begin{enumerate}
\def\labelenumi{\alph{enumi})}
\setcounter{enumi}{1}
\tightlist
\item
  Let \(p\) be a prime number. Let \(Z_{n}\) denote the set \(\mathbf{Z}\) of rational integers, topologized with the inverse image of the discrete topology under the canonical mapping of \(\mathbf{Z}\) onto \(\mathbf{Z} / p^{n}\mathbf{Z}\) . If \(m \leqslant n\) , let \(f_{mn}\) denote the identity mapping \(Z_{n} \to Z_{m}\) , which is continuous. Show that the spaces \(Z_{n}\) are quasi-compact, but that the inverse limit of the inverse system \((Z_{n}, f_{mn})\) is not quasi-compact (*).
\end{enumerate}

\begin{enumerate}
\def\labelenumi{\arabic{enumi})}
\setcounter{enumi}{7}
\tightlist
\item
  Let X be a topological space and let \((\mathrm{G}_{\alpha})_{\alpha\in\mathbb{A}}\) be a family of graphs of equivalence relations on X, whose index set A is directed; let \(R_{\alpha}\) be the relation \((x,y)\in\mathbf{G}_{\alpha}\) and let \(\varphi_{\alpha}\) be the canonical mapping \(X\to X/R_{\alpha}\) . Suppose that the relation \(R_{\beta}\) implies \(R_{\alpha}\) if \(\alpha\leqslant\beta\) and let \(\varphi_{\alpha\beta}\) be the canonical mapping \(X/R_{\beta}\to X/R_{\alpha}\) ; then \((X/R_{\alpha},\varphi_{\alpha\beta})\) is an inverse system of topological spaces. Let \(Y=\varprojlim X/R_{\alpha}\) .
\end{enumerate}

\begin{enumerate}
\def\labelenumi{\alph{enumi})}
\item
  Let \(g: \mathbf{X} \to \mathbf{Y}\) be the continuous mapping \(\varprojlim \varphi_{\alpha}\) . Show that \(g(x)\) is dense in \(\mathbf{Y}\) and that, for each \(x \in \mathbf{X}\) , \(\overrightarrow{g}(g(x))\) is the class of \(x\) with respect to the equivalence relation whose graph is \(\bigcap_{n} G_{\alpha}\) .
\item
  Suppose that all the \(R_{\alpha}\) are closed equivalence relations, and that for each \(x \in X\) and each neighbourhood V of x there is an index \(\alpha\) such that the class of x with respect to \(R_{\alpha}\) is contained in V. Show that under these conditions g is an open mapping of X onto \(g(\mathbf{X})\) .
\item
  Show that if the classes mod \(R_{\alpha}\) are closed and quasi-compact for each \(\alpha \in A\) , then g is surjective.
\end{enumerate}

\begin{enumerate}
\def\labelenumi{\arabic{enumi})}
\setcounter{enumi}{8}
\item
  Let X be an accessible space and \(\mathcal{R}\) a set of quasi-compact closed subsets of X, which contains all finite subsets of X and is such that the union of any two subsets of \(\mathcal{R}\) belongs to \(\mathcal{R}\) . Let \(X'\) be a set which is the sum of X and a set \(\{\omega\}\) consisting of a single point. Show that there exists a quasi-compact accessible topology on \(X'\) which induces the given topology on X, and which is such that the complements in \(X'\) of the sets of \(\mathcal{R}\) form a fundamental system of open neighbourhoods of \(\omega\) . Give examples of distinct topologies on \(X'\) corresponding to distinct sets \(\mathcal{R}\) .
\item
  Let X be a paracompact space. Show that if every open subspace of X is paracompact, then every subspace of X is paracompact. Deduce that every subspace of a locally compact space whose topology has a countable base, is paracompact {[}cf.~Chapter IX, § 4, Exercise 20 a){]}.
\end{enumerate}

¶ 11) Let \(\mathbf{X}_0\) be an uncountable well-ordered set which has a greatest element; let \(a\) be the smallest element of \(\mathbf{X}_0\) and \(b\) the smallest of the elements \(x \in \mathbf{X}_0\) such that the interval \([a, x[\) is uncountable. Let \(\mathbf{X}\) denote the interval \([a, b[\) with the topology \(\mathfrak{C}_{-}(\mathbf{X})\) .

\begin{enumerate}
\def\labelenumi{\alph{enumi})}
\item
  X is locally compact but not compact (Exercise 2) and every sequence in X has a cluster point (observe that every countable subset of X is bounded above).
\item
  Let \(f\) be a mapping of \(\mathbf{X}\) into itself such that \(f(x) < x\) for all sufficiently large \(x\) . Show that there exists an element \(c \in \mathbf{X}\) with the property that for each \(x \in \mathbf{X}\) there exists \(y \geqslant x\) such that \(f(y) \leqslant c\) . {[}Use reductio ad absurdum, by constructing a sequence \((z_n)\) of points of \(\mathbf{X}\) such that \(z_{n+1}\) is the smallest of the elements \(z' \in \mathbf{X}\) such that \(f(x) \geqslant z_n\) whenever \(x \geqslant z'\) {]}.
\item
  Deduce from b) that X is not paracompact.
\end{enumerate}

\begin{enumerate}
\def\labelenumi{\arabic{enumi})}
\setcounter{enumi}{11}
\tightlist
\item
  Let X be a topological space, \(\mathfrak{F}(X)\) the set of non-empty closed subsets of X, \(\mathfrak{K}(X)\) the set of non-empty quasi-compact subsets of X;
\end{enumerate}

endow each of these sets with the topology induced by the topology \(\mathfrak{C}_{\Omega}\) on \(\mathfrak{P}_{0}(\mathbf{X})\) (§ 2, Exercise 7).

\begin{enumerate}
\def\labelenumi{\alph{enumi})}
\tightlist
\item
  Show that if \(\mathfrak{B} \subset \mathfrak{F}(\mathbf{X})\) is quasi-compact with respect to \(\mathcal{C}_{\Omega}\) and if \(\mathbf{X}\) is regular, then the union of the sets \(\mathbf{M} \in \mathfrak{B}\) is closed in \(\mathbf{X}\) .\\
\item
  Show that if \(\mathfrak{B} \subset \mathfrak{K}(\mathbf{X})\) is quasi-compact with respect to \(\mathcal{C}_{\Omega}\) then the union of the sets \(\mathbf{M} \in \mathfrak{B}\) is quasi-compact.
\end{enumerate}

¶ 13) With the notation of Exercise 13, endow \(\mathfrak{F}(\mathbf{X})\) and \(\mathfrak{R}(\mathbf{X})\) with the topology induced by the topology \(\mathfrak{C}_{\Theta}\) on \(\mathfrak{P}_{0}(\mathbf{X})\) (§ 8, Exercise 12). a) Show that if X is quasi-compact, then so is \(\mathfrak{F}(\mathbf{X})\) . (For each open set U in X, let \(c(\mathrm{U})\) {[}resp. \(m(\mathrm{U})\) {]} be the set of all closed subsets M of X such that \(\mathrm{M} \subset \mathrm{U}\) (resp. \(\mathrm{M} \cap \mathrm{U} \neq \emptyset\) ); the \(c(\mathrm{U})\) and \(m(\mathrm{U})\) form a subbase of \(\mathfrak{C}_{\Theta}\) . Then use Exercise 1).

\begin{enumerate}
\def\labelenumi{\alph{enumi})}
\setcounter{enumi}{1}
\item
  Suppose that X is accessible, and let \(X'\) be the image of X under the mapping \(x \to \{x\}\) of X into \(\mathfrak{F}(X)\) . Let H be a subspace of \(\mathfrak{F}(X)\) which contains \(X'\) . Show that if H is quasi-compact then so is X. {[}If \((\mathrm{U}_{\alpha})\) is an open covering of X, remark that if \(B_{\alpha}\) denotes the set of all \(M \in \mathfrak{F}(X)\) such that \(M \cap U_{\alpha} \neq \emptyset\) , then \(\{\mathfrak{B}_{\alpha}\}\) is an open covering of \(\mathfrak{F}(X)]\) .
\item
  Show that if X is locally compact, then \(\mathfrak{K}(\mathbf{X})\) is open in \(\mathfrak{F}(\mathbf{X})\) (use Proposition 10 of § 7).
\item
  Suppose that X is accessible. Show that \(\mathcal{R}(X)\) is Hausdorff (resp. regular, compact, locally compact) if and only if X is. {[}For the first two assertions, use Exercise 12 of § 8 and Exercise 3 of § 9; for the last two assertions, use a) and b) and Exercise 12{]}.
\end{enumerate}

¶ 14) A topological space X is said to be a Lindelöf space if every open covering of X contains a countable covering of X. Every space with a countable base is a Lindelöf space, and so is every quasi-compact space.

\begin{enumerate}
\def\labelenumi{\alph{enumi})}
\item
  Every closed subspace of a Lindelöf space is a Lindelöf space. For every subspace of a Lindelöf space to be a Lindelöf space it is necessary and sufficient that every open subspace is Lindelöf (cf.~Exercise 15).
\item
  If \(f\) is any continuous mapping of a Lindelöf space \(\mathbf{X}\) into a topological space \(\mathbf{X}'\) , then the subspace \(f(\mathbf{X})\) of \(\mathbf{X}'\) is Lindelöf.
\item
  Every countable union of Lindelöf subspaces of a topological space is Lindelöf. In particular every countable space is Lindelöf.
\item
  A Lindelöf space X is quasi-compact if every sequence in X has a cluster point {[}show that axiom (C'') is then verified{]}.
\end{enumerate}

¶ 15) a) Let X be a topological space in which every open subspace is Lindelöf. Show that X satisfies condition (D \(_{III}\) ) of § 1, Exercise 7, and that if X is regular then every closed subset of X is a countable intersection of open sets {[}cf.~Exercise 23 c){]}.

\begin{enumerate}
\def\labelenumi{\alph{enumi})}
\setcounter{enumi}{1}
\item
  Let X be a topological space in which every open subspace is paracompact. Show that every open subspace of X is Lindelöf if and only if X satisfies condition \((D_{\mathrm{III}})\) of §1, Exercise 7.
\item
  Let X be a compact space. Show that every open subspace of X is Lindelöf if and only if (i) X satisfies condition \((\mathbf{D}_{\mathbf{III}})\) of §1, Exercise 7, and (ii) every closed subset of X is a countable intersection of open sets (use b) and Theorem 5 of no. 10).
\end{enumerate}

¶ 16) Let X be a topological space and A a subset of X. A point \(a \in X\) is said to be a condensation point of A if every neighbourhood of \(a\) contains an uncountable infinity of points of A.

\begin{enumerate}
\def\labelenumi{\alph{enumi})}
\item
  The set of condensation points of a subset A of X is closed. Give an example in which X has only one condensation point (cf.~no. 8, Theorem 4).
\item
  Suppose X is an accessible space in which every open subspace is Lindelöf (Exercise 15). Show that for every subset A of X, the set B of condensation points of A is perfect, and that A ∩ GB is countable.
\end{enumerate}

¶ 17) In a topological space X, a point \(a\) is said to be totally adherent to a subset A of X if for each neighbourhood U of \(a\) we have card (A) = card (A ∩ U). Show that X is quasi-compact if and only if for each infinite subset A of X there is a point of X which is totally adherent to A. (If X is not quasi-compact, show that there is an initial ordinal \(\omega_{\alpha}\) (Set Theory, Chapter III, § 6, Exercise 10) and an open covering (U \(_{\xi}\) ) of X, where \(\xi\) runs through the set of ordinals \(<\omega_{\alpha}\) , such that for each index \(\xi<\omega_{\alpha}\) the complement A \(_{\xi}\) in X of the union of the U \(_{\eta}\) with indices \(\eta<\xi\) has a cardinal equal to N \(_{\alpha}\) ; deduce that it is possible to define a family ( \(x_{\xi}\) ) ( \(\xi<\omega_{\alpha}\) ) of points of X by transfinite induction such that \(x_{\xi}\in\mathrm{A}_{\xi}\) for each \(\xi\) and \(x_{\xi}\neq x_{\eta}\) whenever \(<\eta\xi\) ).

¶ 18) A Hausdorff space X is said to be absolutely closed if it has the following property: for every homeomorphism f of X onto a subspace of a Hausdorff space X', f(X) is closed in X'. Every compact space is absolutely closed.

\begin{enumerate}
\def\labelenumi{\alph{enumi})}
\tightlist
\item
  Show that the following conditions are equivalent:
\end{enumerate}

(AF) The space X is absolutely closed.

(AF') Every filter base on X which consists of open sets has at least one cluster point.

(AF'') Every family of closed subsets of X whose intersection is empty contains a finite subfamily such that the interiors of the subsets in this subfamily have an empty intersection.

(AF''') Every open covering of X contains a finite subfamily whose closures cover X.

\begin{enumerate}
\def\labelenumi{\alph{enumi})}
\setcounter{enumi}{1}
\item
  Let X be a Hausdorff space and U a dense open subset of X. Suppose that every filter base on U which consists of open sets has at least one cluster point in X. Show that under these conditions X is absolutely closed. In particular, if A is an open subset of an absolutely closed space X, then \(\overline{A}\) is absolutely closed.
\item
  Let X be an absolutely closed space and let f be a continuous mapping of X into a Hausdorff space \(X'\) . Show that \(f(\mathbf{X})\) is an absolutely closed subspace of \(X'\) .
\item
  Show that the product of two absolutely closed spaces is absolutely closed (argue as in Proposition 17 of no. 10).
\end{enumerate}

¶ 19) A Hausdorff space X is said to be minimal if its topology C is such that every topology on X which is Hausdorff and coarser than C necessarily coincides with C; then C must be semi-regular (§ 8, Exercise 20). Show that a Hausdorff space X is minimal if and only if every filter base on X, which consists of open sets and has at most one cluster point, is convergent. It comes to the same thing to say that X is absolutely closed (Exercise 18) and that every filter base on X, which consists of open sets and has a single cluster point, converges to this point.

¶ 20) A topological space X is said to be completely Hausdorff (and the topology of X is said to be completely Hausdorff) if each pair of distinct points of X has disjoint closed neighbourhoods.

\begin{enumerate}
\def\labelenumi{\alph{enumi})}
\item
  Every regular space is completely Hausdorff. Every subspace of a completely Hausdorff space is completely Hausdorff. Every product of completely Hausdorff spaces is completely Hausdorff. Every topology which is finer than a completely Hausdorff topology is completely Hausdorff.
\item
  Every minimal completely Hausdorff space X (Exercise 19) is compact. (By reductio ad absurdum: consider an ultrafilter U on X and the filter base formed by the open sets which belong to U).
\end{enumerate}

\(^{*}c)\) Let \(\theta\) be an irrational number \textgreater0. For every point \((x,y)\in\mathbf{Q}^{2}\) and every \(\alpha>0\) let \(\mathrm{B}_{\alpha}(x,y)\) be the set consisting of \((x,y)\) and the points \((z,0)\in\mathbf{Q}^{2}\) such that \(|z-(x+\theta y)|<\alpha\) or \(|z-(x-\theta y)|<\alpha\) . Let \(\mathfrak{B}(x,y)\) be the set of \(\mathrm{B}_{\alpha}(x,y)\) where \(\alpha\) runs through the set of real numbers \textgreater0. Show that there is a Hausdorff topology \(\mathcal{C}\) on \(\mathbf{Q}^{2}\) such that, for each \((x,y)\) , \(\mathfrak{B}(x,y)\) is a fundamental system of neighbourhoods of \((x,y)\) for \(\mathcal{C}\) . Show that, given any two points a, b of \(\mathbf{Q}^{2}\) , every closed neighbourhood of a meets every closed neighbourhood of b in the topology \(\mathcal{C}.*\)

¶ 21) a) Let X be a Hausdorff space and let \(\mathfrak{G}\) be the topology of X. Show that X is absolutely closed (Exercise 18) if and only if the semi-regular topology \(\mathfrak{G}^*\) associated with \(\mathfrak{G}\) ( \(\S\) 8, Exercise 20) is the topology of a minimal space (Exercise 19).

\begin{enumerate}
\def\labelenumi{\alph{enumi})}
\setcounter{enumi}{1}
\tightlist
\item
  If \(\mathfrak{C}\) is completely Hausdorff (Exercise 20), then so is \(\mathfrak{C}_{*}\) (use Exercise 20 of § 8). Deduce that an absolutely closed regular space is compact.
\end{enumerate}

¶ 22) a) Let X be a Hausdorff space and let A be a dense subset of X. Show that if every filter on A has at least one cluster point in X, then X is absolutely closed {[}verify axiom (AF') of Exercise 18{]}.

\begin{enumerate}
\def\labelenumi{\alph{enumi})}
\setcounter{enumi}{1}
\tightlist
\item
  Let \(X_{0}\) be a compact space which has a non-open dense subset A, and let \(C_{0}\) be the topology of \(X_{0}\) . Let C be the topology on \(X_{0}\) generated by A and the subsets of \(X_{0}\) which are open in \(C_{0}\) , so that \(C^{*} = C_{0}\) (§ 8, Exercise 20). If X denotes the non-compact absolutely closed space obtained by endowing the set \(X_{0}\) with the topology C (Exercise 21) show that every filter on A has a cluster point in X.
\end{enumerate}

¶ 23) a) Let X be a countable absolutely closed space. Show that the set of isolated points of X is dense in X. {[}Use reductio ad absurdum, by arranging the points of X in a sequence, and constructing by induction a countable open covering (U \(_{n}\) ) of X such that no finite union of the \(\overline{U}_{n}\) covers X{]}.

*b) Give an example of a completely Hausdorff absolutely closed topological space which is not a Lindelöf space {[}in the product \([0,1]\times[0,1]\) ; consider the complements of the sets \(A\times\{o\}\) , where A runs through the set of all subsets of \([0,1]\) , and apply the general method of § 8, Exercise 20 c){]}. This space satisfies condition \((D_{II})\) of § 1, Exercise 1, but does not satisfy condition \((D_{III})\) .

\begin{enumerate}
\def\labelenumi{\alph{enumi})}
\setcounter{enumi}{2}
\tightlist
\item
  Let \(X_{0}\) be a compact space with no isolated points, let M be the set of complements of countable subsets of \(X_{0}\) {[}so that the sets of M are dense in \(X_{0}\) , cf.~a){]}. Let C be the topology generated by the open sets of \(X_{0}\) and the sets of M. Show that the space X obtained by giving \(X_{0}\) the topology C is absolutely closed and completely Hausdorff and is a Lindelöf space which does not satisfy condition ( \(D_{II}\) ) of §1, Exercise 7; show also that no point of X has a countable fundamental system of neighbourhoods and that no sequence in X, all of whose terms are distinct, has a cluster point. If every open subspace of \(X_{0}\) is Lindelöf {[}cf.~Exercise 16 c){]}, show that the same is true of X and hence in particular that X satisfies ( \(D_{III}\) ) {[}Exercise 16 a){]}. This is the case when \(X_{0} = [0,1]\) . Show however that in this case there are countable closed subsets of X which are not countable intersections of open sets of X (cf.~Chapter IX, §5, no. 3). *
\end{enumerate}

¶ 24) a) Let X be a Hausdorff space, and let Γ be the set of filters on X which have a base consisting of open sets and which have no cluster point in X. Show that the set Γ is inductive if it is not empty (i.e., if X is not absolutely closed).

\begin{enumerate}
\def\labelenumi{\alph{enumi})}
\setcounter{enumi}{1}
\tightlist
\item
  Let \(\Omega\) be the set of maximal filters in \(\Gamma\) , and let \(X'\) be the set which is the sum of \(X\) and \(\Omega\) . For each \(x \in X\) let \(\mathfrak{B}(x)\) denote the set of neighbourhoods of \(x\) in \(X\) , and for each \(\omega \in \Omega\) let \(\mathfrak{B}(\omega)\) denote the set of subsets of \(X'\) of the form \(\{\omega\} \cup M\) , where \(M\) belongs to the filter \(\omega\) . Then there is a topology on \(X'\) in which \(\mathfrak{B}(x')\) is a fundamental system of neighbourhoods of \(x'\) for each \(x' \in X'\) . Show that, with this topology, \(X'\) is absolutely closed and has the following universal property: given any open continuous mapping \(f: X \to Y\) of \(X\) into an absolutely closed space \(Y\) , there is a unique continuous mapping \(g: X' \to Y\) which extends \(f\) (consider the images under \(f\) of the filters which belong to \(\Omega\) , and observe that in an absolutely closed space, a filter which is maximal in the set of filters which have a base consisting of open sets must be convergent). In particular, the topology of \(X'\) is X-maximal ( \(\S 3\) , Exercise 11). Deduce that there are absolutely closed spaces whose topology is quasi-maximal ( \(\S 2\) , Exercise 6).
\end{enumerate}

¶ 25) a) Let X be an accessible space, and let \(\Phi\) be the set of filters on X which have a base consisting of closed sets. Show that the set \(\Phi\) is inductive, and let \(X''\) be the set of maximal filters in \(\Phi\) .

\begin{enumerate}
\def\labelenumi{\alph{enumi})}
\setcounter{enumi}{1}
\tightlist
\item
  For each open subset U of X, let \(U^{*}\) be the set of filters \(F \in X''\) such that \(U \in F\) . If U, V are two open subsets of X, then
\end{enumerate}

\[
(\mathrm{U} \cap \mathrm{V}) ^ {*} = \mathrm{U} ^ {*} \cap \mathrm{V} ^ {*} \quad \text { and } \quad (\mathrm{U} \cup \mathrm{V}) ^ {*} = \mathrm{U} ^ {*} \cup \mathrm{V} ^ {*}
\]

(remark that if \(\mathfrak{F} \in \mathbf{X}''\) is such that \(\mathbf{U} \notin \mathfrak{F}\) , then \(\mathbf{X} = \mathbf{U} \in \mathfrak{F}\) ). Show that the sets \(\mathbf{U}^*\) form a base of a topology on \(\mathbf{X}''\) and that in this topology \(\mathbf{X}''\) is quasi-compact {[}verify axiom \((\mathbf{C}''')\) , using the remark above{]} and accessible {[}if \(\mathfrak{F}, \mathfrak{G}\) are two distinct elements of \(\mathbf{X}''\) , show that there exist two disjoint closed subsets \(\mathbf{A}\) and \(\mathbf{B}\) in \(\mathbf{X}\) such that \(\mathbf{A} \in \mathfrak{F}\) and \(\mathbf{B} \in \mathfrak{G}\) {]}.

\begin{enumerate}
\def\labelenumi{\alph{enumi})}
\setcounter{enumi}{2}
\item
  For each \(x \in \mathbf{X}\) let \(\varphi(x)\) be the ultrafilter on \(\mathbf{X}\) which has the single set \(\{x\}\) as a base. Show that \(\varphi\) is a homeomorphism of \(\mathbf{X}\) onto a dense subset of \(\mathbf{X}''\) and that if \(\mathbf{X}\) is quasi-compact (and accessible), then \(\varphi(\mathbf{X}) = \mathbf{X}''\) .
\item
  Show that for each continuous mapping \(f\) of \(\mathbf{X}\) into a compact space \(\mathbf{Y}\) there is a unique continuous mapping \(g: \mathbf{X}'' \to \mathbf{Y}\) such that \(f = g \circ \varphi\) .
\item
  Let \(\tilde{\mathbf{X}}\) be the universal Hausdorff space associated with \(\mathbf{X}''\) (§ 8, Exercise 27); show that the canonical mapping \(\psi: \mathbf{X}'' \to \tilde{\mathbf{X}}\) is surjective and that \(\tilde{\mathbf{X}}\) is compact. Thus every continuous mapping of \(\mathbf{X}\) into a compact space \(\mathbf{Y}\) can be written uniquely in the form \(f = h \circ \psi \circ \varphi\) , where \(h: \tilde{\mathbf{X}} \to \mathbf{Y}\) is continuous.
\end{enumerate}

\begin{enumerate}
\def\labelenumi{\arabic{enumi})}
\setcounter{enumi}{25}
\tightlist
\item
  If X is a discrete space, show that the spaces \(X''\) and \(\tilde{X}\) defined in Exercise 25 are identical, and that if C is the topology of the space \(X'\) defined in Exercise 24, then \(X''\) can be identified with the space \(X'\) carrying the semi-regular topology \(C^*\) associated with C (§ 8, Exercise 20). Furthermore, the set \(X''\) can be canonically identified with the set of ultrafilters on X. The compact space \(X''\) is called the ultrafilter space of X.
\end{enumerate}

¶ 27) Show that in a compact space X every isodyne {[}§ 2, Exercise 11 d){]} infinite open subspace is solvable. {[}Using Proposition 1 of no. 2, show that there is a base B of the topology of U such that card (B) = card (U); then use Exercise 11 c) of § 2.{]} Deduce that every locally compact space with no isolated points is solvable {[}use Exercise 11 d) of § 2{]}; consequently such a space is not submaximal.

\begin{enumerate}
\def\labelenumi{\arabic{enumi})}
\setcounter{enumi}{27}
\tightlist
\item
  \begin{enumerate}
  \def\labelenumii{\alph{enumii})}
  \tightlist
  \item
    Let K be the discrete space \(\{0,1\}\) . If X is any compact space, let F be the set of mappings f of K into the set \(\mathfrak{F}(X)\) of closed subsets of X such that \(\mathbf{X}=f(0)\cup f(1)\) . Consider the compact space \(Y=K^{F}\) ; show that for each \(y=(y_{f})_{f\in\mathbb{F}}\) in Y, the intersection of the closed sets \(f(y_{f})\) of X consists of at most one point. The set Z of \(y\in Y\) for which this intersection is not empty is a closed subset of Y. If \(\varphi(y)\) is the unique point common to the \(f(y_{f})\) for \(y\in Z\) , show that \(\varphi\) is a continuous surjection of Z onto X.
  \end{enumerate}
\end{enumerate}

\begin{enumerate}
\def\labelenumi{\alph{enumi})}
\setcounter{enumi}{1}
\tightlist
\item
  Give an example of a compact space X such that there is no surjective continuous mapping of a compact space of the form \(K^{A}\) onto X {[}use Exercise 4 b) of § 4 and Exercise 26 of § 9{]}.
\end{enumerate}

§ 29) A topological space X is said to be locally quasi-compact if each point of X has a fundamental system of quasi-compact neighbourhoods. a) Show that in a locally quasi-compact space, every quasi-compact subset has a fundamental system of quasi-compact neighbourhoods.

\begin{enumerate}
\def\labelenumi{\alph{enumi})}
\setcounter{enumi}{1}
\tightlist
\item
  Show that every locally closed subspace of a locally quasi-compact space is locally quasi-compact.
\end{enumerate}

\(^{*}c)\) On the closed Euclidean disc B: \(||x|| \leqslant 1\) in the real plane \(R^{2}\) , let \(\mathcal{G}_{0}\) be the topology induced by that of \(R^{2}\) , and let \(\mathcal{G}\) be the (finer) topology defined as follows: the neighbourhood filter of a point \(x \neq (1,0)\) of B is the same for \(\mathcal{G}\) as for \(\mathcal{G}_{0}\) , but for the point \((1,0)\) a fundamental system of neighbourhoods in the topology \(\mathcal{G}\) is formed by the sets

\[
\left\{\left(\mathrm{I}, \mathrm{o}\right) \right\} \cup \left(\mathbf {D} \cap \mathbf {B} _ {\mathbf {0}}\right),
\]

where D is any open disc centred at (1,0) and \(B_{0}\) is the open disc \(\|x\|<1\) . Let Y be the set B with the topology C, and let R be the equivalence relation on Y whose equivalence classes are the set S of all \(x\neq(1,0)\)

in B such that \(\|x\| = 1\) and the subsets of Y which consist of a single point of B - S. Show that the quotient space X = Y/R is quasi-compact but not locally quasi-compact. *

